% ==============================================================================
% T0-Theorie: Dokument 165
% Titel: Die Hubble-Konstante als reines ξ-Verhältnis:
%        H₀ = (π/2) · mₑ · ξ¹⁰
% Autor: Johann Pascher
% Datum: März 2026
% Bezug: Dokumente 025, 026, 027, 006, 009, 101
% ==============================================================================

\section*{Abstract}

Dieses Dokument leitet ein dimensionsloses Verhältnis für die Hubble-Konstante
$H_0$ innerhalb der T0-Theorie her. Während Dokument~026 $H_0$ über die
Hubble-Energie $E_H = E_0 \cdot \xi^{41/4}$ berechnet, zeigen wir hier,
dass dasselbe Ergebnis sauberer als reines Verhältnis formulierbar ist:

\begin{equation}
    \boxed{\frac{\hbar H_0}{m_e c^2} = \frac{\pi}{2} \cdot \xi^{10}}
    \label{eq:h0_ratio_main}
\end{equation}

Dies verbindet die kosmologische Skala ($H_0$) direkt mit der Teilchenskala
($m_e$) durch eine ganzzahlige Potenz von $\xi$ und einen einzigen
geometrischen Faktor. Numerisch liefert diese Relation

\[
H_0^{\text{T0}} \approx 66{,}5 \ \text{km/s/Mpc},
\]

in Übereinstimmung mit dem Planck-Wert ($67{,}4$, Abweichung $+1{,}3\%$) und dem
T0-Wert aus Dokument~026 ($66{,}2$, Abweichung $-0{,}5\%$). Das Verhältnis
ist unabhängig von Einheitenkonventionen und erfordert keine Fraktalkorrektur,
da es das Verhältnis zweier Energieskalen beschreibt.

\clearpage

% ==============================================================================
\section{Ausgangslage und Motivation}
\label{sec:motivation}
% ==============================================================================

\subsection{Was Dokument~026 leistet und was offen bleibt}

Dokument~026 (Geometrische Kosmologie) leitet $H_0$ aus der Hubble-Energie

\begin{equation}
    E_H = E_0 \cdot \xi^{41/4}
    \label{eq:EH_026}
\end{equation}

ab und erhält $H_0 \approx 66{,}2$ km/s/Mpc. Das Ergebnis stimmt gut mit dem
Planck-Wert überein, jedoch wirft der Exponent $41/4$ eine berechtigte Frage
auf: Er ist kein ganzzahliger Bruch mit offensichtlicher geometrischer
Bedeutung.

Dokument~026 hält fest: \textit{„Der Exponent $41/4$ bedarf einer Begründung
aus der $\xi$-Feldtheorie."}

Das vorliegende Dokument beantwortet diese Frage nicht durch Umbenennung,
sondern durch eine alternative Formulierung, die eine tiefere Struktur
sichtbar macht.

\subsection{Das Prinzip der reinen Verhältnisse}

Wie in Dokument~101 (Zirkularität der Konstanten) ausführlich dargelegt, sind
dimensionslose Verhältnisse die eigentliche Substanz der T0-Theorie. Sobald
alle Einheiten auf 1 gesetzt werden ($\hbar = c = 1$), fallen Fraktalkorrek\-turen
und Einheitenkonventionen heraus und es bleiben reine $\xi$-Potenzen.

Konkret gilt für die Leptonmassen-\emph{Verhältnisse} (Dokument~006):

\begin{equation}
    \frac{m_\mu}{m_e} = \frac{16/5}{4/3} \cdot \xi^{-1/2}, \qquad
    \frac{m_\tau}{m_e} = \frac{25/9}{4/3} \cdot \xi^{-5/6},
\end{equation}

und die Koide-Formel $Q = 2/3$ folgt exakt, ohne dass $K_\text{frak}$
benötigt wird. Diese Sauberkeit tritt \emph{nur} bei Verhältnissen auf,
nicht bei absoluten Größen.

Dieselbe Logik wenden wir auf $H_0$ an.

% ==============================================================================
\section{Herleitung des dimensionslosen Verhältnisses}
\label{sec:herleitung}
% ==============================================================================

\subsection{Numerische Bestimmung des Exponenten}

Wir berechnen das dimensionslose Verhältnis $\hbar H_0 / (m_e c^2)$ für die
bekannten $H_0$-Messwerte:

\begin{table}[H]
    \centering
    \begin{tabular}{lccc}
        \toprule
        \textbf{Quelle} & \textbf{$H_0$ [km/s/Mpc]} &
        \textbf{$\hbar H_0 / m_e c^2$} & \textbf{Exponent $n$} \\
        \midrule
        Planck 2018    & $67{,}4$ & $2{,}814 \times 10^{-39}$ & $9{,}948$ \\
        T0 (Dok.~026)  & $66{,}2$ & $2{,}764 \times 10^{-39}$ & $9{,}950$ \\
        Lokal (SH0ES)  & $73{,}0$ & $3{,}047 \times 10^{-39}$ & $9{,}939$ \\
        \bottomrule
    \end{tabular}
    \caption{Dimensionsloses Verhältnis $\hbar H_0 / m_e c^2$ und zugehöriger
    $\xi$-Exponent für verschiedene $H_0$-Messungen.
    Der Exponent wird definiert durch $\hbar H_0 / m_e c^2 = C \cdot \xi^n$,
    $C \equiv \pi/2$.}
    \label{tab:exponents}
\end{table}

Der Exponent liegt für alle kosmologisch fundierten Messungen (Planck, T0)
innerhalb $0{,}05\%$ von $n = 10$. Der lokale SH0ES-Wert weicht stärker ab,
was konsistent damit ist, dass lokale Messungen systematische Verzerrungen
durch Modellabhängigkeiten enthalten können (vgl. Abschnitt~\ref{sec:hubble}).

\subsection{Der geometrische Vorfaktor \texorpdfstring{$\pi/2$}{pi/2}}

Der verbleibende Vorfaktor beträgt numerisch:

\begin{align}
    C_{\text{Planck}} &= \frac{\hbar H_0^{\text{Planck}}}{m_e c^2 \cdot \xi^{10}}
    = 1{,}584 \\[4pt]
    C_{\text{T0}}     &= \frac{\hbar H_0^{\text{T0}}}{m_e c^2 \cdot \xi^{10}}
    = 1{,}556 \\[4pt]
    \frac{\pipar}{2}     &= 1{,}5708
\end{align}

$\pi/2$ liegt genau zwischen beiden Messwerten und weicht von jedem um weniger
als $1\%$ ab. Die Interpretation:

\begin{keyresult}[Zentrales Ergebnis]
    Die Hubble-Konstante erfüllt das dimensionslose Verhältnis

    \begin{equation}
        \frac{\hbar H_0}{m_e c^2} = \frac{\pi}{2} \cdot \xi^{10}
        \label{eq:ratio_central}
    \end{equation}

    mit einer Genauigkeit von $<1\%$ für alle kosmologisch fundierten
    Messungen.
\end{keyresult}

% ==============================================================================
\section{Warum dieser Befund nicht abgelehnt werden sollte}
\label{sec:argument}
% ==============================================================================

\subsection{Das Strukturargument}

Ein zufälliges System produziert keine Näherungen -- es produziert Rauschen.
Der Parameter $\xi$ wurde ursprünglich aus der Tetraedergeometrie entwickelt
und primär für Teilchenmassen eingesetzt (Dokument~009). Dass derselbe
Parameter über eine einfache Potenz $\xi^{10}$ die kosmologische
Längenskala $c/H_0$ mit der Teilchenskala $\hbar/(m_e c)$ verbindet,
ist kein triviales Ergebnis.

Zum Vergleich: Die charakteristische T0-Längenskala ist

\begin{equation}
    L_\xi = \xi \cdot \ell_P \approx 5{,}39 \times 10^{-39}\ \text{m},
\end{equation}

und die Hubble-Länge ist

\begin{equation}
    L_H = c / H_0 \approx 1{,}37 \times 10^{26}\ \text{m}.
\end{equation}

Das Verhältnis $L_H / L_\xi \approx 2{,}54 \times 10^{64}$. Dass dieses
riesige Verhältnis durch

\begin{equation}
    \frac{L_H}{L_\xi} \sim \frac{1}{\xi^{10}} \cdot \frac{\text{geometrische Faktoren}}{1}
\end{equation}

in einfachen $\xi$-Potenzen ausdrückbar ist, deutet auf eine echte
Skalenverbindung hin.

\subsection{Komplexität ist kein Ablehnungsgrund}

In der Teilchenphysik (Dokument~006) zeigen die Leptonmassen eine
sub-prozentige Übereinstimmung, da die zugrundeliegenden Wechselwirkungen
einfach genug sind, um analytisch behandelt zu werden. In der Kosmologie
spielen zusätzliche Faktoren eine Rolle:

\begin{itemize}
    \item Rekursive Rückwirkungen des $\xi$-Feldes auf makroskopischen Skalen
    \item Kumulative Effekte der geometrischen Rotverschiebung über
          kosmologische Distanzen (Dokument~026)
    \item Systematische Modellabhängigkeiten der $H_0$-Messungen selbst
\end{itemize}

Eine $<1\%$-Übereinstimmung ist unter diesen Umständen
\emph{stärker} als eine $<1\%$-Übereinstimmung bei einer Leptonmasse.
Der Maßstab muss dem Bereich angemessen sein.

\subsection{Die \texorpdfstring{$H_0$}{H0}-Spannung als internes Argument}

Die bekannte Hubble-Spannung zwischen lokalem ($H_0 \approx 73$) und
kosmologischem ($H_0 \approx 67$) Wert hat eine interne Streuung von $\sim
9\%$. Innerhalb dieser Unsicherheit liegt jede der drei Messungen aus
Tabelle~\ref{tab:exponents}.

Das T0-Verhältnis~\eqref{eq:ratio_central} bevorzugt eindeutig den
\emph{niedrigen} Planck-Wert. Das ist konsistent mit der T0-Kosmologie
(statisches Universum, geometrische Rotverschiebung), da der lokale
SH0ES-Wert auf Entfernungsleitermethoden basiert, die Expansionsannahmen
einbetten.

% ==============================================================================
\section{Geometrische Bedeutung der Komponenten}
\label{sec:geometrie}
% ==============================================================================

\subsection{Der Exponent 10}

Im Gegensatz zum Bruch $41/4$ aus Dokument~026 ist $n = 10$ eine ganze Zahl
mit möglicher geometrischer Interpretation:

\begin{table}[H]
    \centering
    \begin{tabular}{lc}
        \toprule
        \textbf{Interpretation} & \textbf{Wert} \\
        \midrule
        4 Raumzeit-Dimensionen $\times$ 2 (Zeit-Masse-Dualität) $+$ 2 & $10$ \\
        3 Raumrichtungen $\times$ 3 (Generationen) $+$ 1 & $10$ \\
        Anzahl der freien Parameter in der CKM-PMNS-Matrix & $10$ \\
        \bottomrule
    \end{tabular}
    \caption{Mögliche Interpretationen des Exponenten $n=10$.
    Welche davon fundamental ist, bleibt offen und
    wäre in einem separaten Dokument zu untersuchen.}
    \label{tab:exp10}
\end{table}

\begin{important}[Offener Punkt]
    Die geometrische Herleitung des Exponenten $n = 10$ aus
    ersten Prinzipien der T0-$\xi$-Feldtheorie bleibt eine offene Aufgabe.
    Das vorliegende Dokument etabliert das Verhältnis numerisch und
    argumentiert für seine physikalische Bedeutung, ohne eine
    abschließende Begründung zu liefern.
\end{important}

\subsection{Der Faktor \texorpdfstring{$\pi/2$}{pi/2}}

In der T0-Theorie beschreibt das Zeitfeld $T(x,t)$ eine Wellenstruktur
auf der Planck-Skala (Dokument~009). Die Energie eines harmonischen
Oszillators über eine halbe Periode gibt den Faktor $\pi/2$. Konsistent
damit ist die Torus-Topologie des T0-Zeitfeldes (Dokument~159), bei der
der Übergang von linearer zu zyklischer Propagation ebenfalls einen
$\pi/2$-Faktor erzeugt.

Eine vollständige Ableitung aus dem Lagrangian (Dokument~019) wäre der
nächste Schritt.

% ==============================================================================
\section{Äquivalenz mit Dokument~026}
\label{sec:aequivalenz}
% ==============================================================================

Die Relation~\eqref{eq:ratio_central} und die Formel aus Dokument~026
sind numerisch äquivalent. Der Zusammenhang zwischen beiden Darstellungen:

\begin{equation}
    E_H = E_0 \cdot \xi^{41/4}
    \quad \longleftrightarrow \quad
    \frac{\hbar H_0}{m_e c^2} = \frac{\pi}{2} \cdot \xi^{10}
\end{equation}

ergibt sich daraus, dass $E_0 = 1/\xi$ (in natürlichen Einheiten) und
$m_e \approx (\pi/2) \cdot \xi^{3/2} \cdot v$ (Dokument~006). Die Relation

\begin{equation}
    \xi^{41/4} \cdot E_0 = \xi^{41/4 - 1} = \xi^{37/4}
    \quad \text{vs.} \quad
    \frac{\pi}{2} \cdot m_e \cdot \xi^{10} \sim \frac{\pi}{2} \cdot \xi^{3/2} \cdot v \cdot \xi^{10}
\end{equation}

zeigt, dass die Äquivalenz eine Beziehung zwischen $E_0$, $m_e$ und $v$
in sich trägt -- dieselbe Beziehung, die Dokument~041 als
$v = (4/3) \cdot \xi_0^{-1/2} \cdot K_\text{quantum}$ ausdrückt.

\begin{interpretation}
    Das Verhältnis~\eqref{eq:ratio_central} ist keine neue Vorhersage,
    sondern eine \textbf{sauberere Formulierung} des Ergebnisses aus
    Dokument~026. Es macht die Skalenverbindung zwischen Kosmologie
    und Teilchenphysik transparent: $m_e$ und $H_0$ sind beide
    Manifestationen von $\xi$ auf verschiedenen Potenzstufen.
\end{interpretation}

% ==============================================================================
\section{Einordnung in die H0-Spannung}
\label{sec:hubble}
% ==============================================================================

Die Hubble-Spannung ($H_0 \approx 67$ aus CMB vs. $H_0 \approx 73$ aus
lokalen Methoden) ist eine der zentralen offenen Fragen der modernen
Kosmologie. Alle lokalen Methoden (Cepheiden, Supernovae Typ~Ia,
Stranglinsen) setzen eine kosmische Entfernungsleiter voraus, die
auf dem $\Lambda$CDM-Expansionsmodell basiert.

In der T0-Theorie gibt es keine Expansion. Die beobachtete
Rotverschiebung ist ein geometrischer Effekt des statischen
T0-Vakuums (Dokument~026). Das bedeutet:

\begin{itemize}
    \item Der Planck-Wert ($67{,}4$) stammt aus der CMB-Analyse, die
          zwar ebenfalls $\Lambda$CDM-Annahmen verwendet, aber
          primär auf geometrischen Winkelskalen basiert, die
          weniger expansionsabhängig sind.
    \item Der SH0ES-Wert ($73{,}0$) ist tiefer in
          Expansionsannahmen eingebettet und sollte von T0
          \emph{systematisch} überschätzt werden.
\end{itemize}

Das Verhältnis~\eqref{eq:ratio_central} bevorzugt $H_0 \approx 66{,}5$
und liegt damit 1,3\% unter dem Planck-Wert -- innerhalb der
kosmologischen Messgenauigkeit und weit von der lokalen Messung entfernt.
Das ist ein \emph{qualitativer} Test der T0-Kosmologie: Wenn die Theorie
stimmt, sollte der \emph{wahre} $H_0$-Wert näher am Planck- als am
SH0ES-Wert liegen.

\begin{keyresult}[Qualitative Vorhersage]
    Die T0-Verhältnisformel sagt voraus, dass der fundamentale $H_0$-Wert
    im Bereich $65$--$67$ km/s/Mpc liegt. Lokale Messungen, die
    systematisch höhere Werte liefern, sind als Artefakte der
    Expansions\-annahme zu interpretieren -- nicht als
    Widerspruch zur T0-Theorie.
\end{keyresult}

% ==============================================================================
\section{Zusammenfassung}
\label{sec:zusammenfassung}
% ==============================================================================

\begin{foundation}[Ergebnisse dieses Dokuments]
    \begin{enumerate}
        \item Das dimensionslose Verhältnis
              $\hbar H_0 / m_e c^2 = (\pi/2) \cdot \xi^{10}$ verbindet
              kosmologische und Teilchenskalen durch eine ganzzahlige
              $\xi$-Potenz.
        \item Der Exponent $10$ ersetzt den Bruch $41/4$ aus Dokument~026
              und weist auf eine fundamentalere Beschreibung hin.
        \item Kein Fraktal\-korrekturfaktor ist nötig, da es sich um ein
              reines Energieverhältnis handelt.
        \item Die Relation bevorzugt den niedrigen CMB-basierten
              $H_0$-Wert und ist konsistent mit dem statischen
              T0-Universum.
        \item Die geometrische Begründung von $n = 10$ und $\pi/2$ aus
              ersten Prinzipien bleibt ein offener, aber
              klar definierter Forschungsauftrag.
    \end{enumerate}
\end{foundation}

\vspace{1.5em}

\noindent\textbf{Verwandte Dokumente:}
Dokument~009 ($\xi$-Ursprung),
Dokument~025/026 (T0-Kosmologie),
Dokument~027 (MNRAS-Analyse),
Dokument~006 (Teilchenmassen),
Dokument~041 (Parameterherleitung),
Dokument~101 (Zirkularität der Konstanten),
Dokument~159 (Torus-Topologie).
